%==============================================================================
% Standalone NTC Engagement Brief
% A short, printable handout for the meeting with NTC Pokhara engineers.
% Compile separately: pdflatex ntc_engagement_brief.tex
%==============================================================================
\documentclass[11pt,a4paper]{article}
\usepackage[utf8]{inputenc}
\usepackage[T1]{fontenc}
\usepackage{lmodern}
\usepackage[a4paper,left=1in,right=1in,top=0.9in,bottom=0.9in]{geometry}
\usepackage{setspace}
\onehalfspacing
\usepackage{graphicx}
\usepackage{enumitem}
\usepackage{titlesec}
\usepackage{booktabs}
\usepackage{tabularx}
\usepackage{hyperref}
\hypersetup{colorlinks=false}

\titleformat{\section}{\large\bfseries}{\thesection.}{0.5em}{}
\titleformat{\subsection}{\normalsize\bfseries}{\thesubsection.}{0.4em}{}

\begin{document}

\begin{center}
    {\Large\bfseries Engagement Brief}\\[0.2cm]
    {\large SON-GNN-DQN Major Project, Paschimanchal Campus (IOE,
    Tribhuvan University)}\\[0.2cm]
    {\small Prepared for the meeting with Nepal Telecom (NTC) Pokhara
    LTE Optimization Team}
\end{center}

\section{Who we are}

We are four final-year students of the Bachelor of Electronics,
Communication and Information Engineering at Paschimanchal Campus,
Institute of Engineering, Tribhuvan University. Our major project is
supervised by Asst.\ Prof.\ Er.\ Hom Nath Tiwari, Head of the Department
of Electronics and Computer Engineering.

\begin{center}
\begin{tabular}{ll}
Rojin Puri      & PAS078BEI028 \\
Rujan Subedi    & PAS078BEI030 \\
Saange Tamang   & PAS078BEI031 \\
Sulav Kandel    & PAS078BEI041 \\
\end{tabular}
\end{center}

\section{What the project is about (one-paragraph summary)}

The project develops \textbf{SON-GNN-DQN}, a topology-generalized Graph
Neural Network and Deep Q-Learning framework for LTE handover and load
optimization, wrapped in a safety-bounded Self-Organizing Network (SON)
translation layer. The GNN encodes the cellular topology as a dynamic
graph; the Deep Q-Network selects the best target cell for a UE; the
SON layer converts the learned preference into bounded Cell Individual
Offset (CIO) and Time-to-Trigger (TTT) updates that respect 3GPP and
operator-side constraints. The framework is designed to be deployable
on top of existing infrastructure (no replacement of vendor equipment),
and compatible with a future O-RAN xApp deployment.

\section{What we have built so far}

\begin{itemize}[leftmargin=*]
    \item A complete LTE mobility simulator in Python (RSRP, RSRQ, SINR,
    throughput, PRB load, handover events).
    \item A GNN encoder (Graph Attention Network) and a Double + Dueling
    Deep Q-Network agent with prioritized replay.
    \item Six baseline mobility policies (no-handover, random-valid,
    strongest-RSRP, A3+TTT, load-aware, and our learned policy with and
    without the SON translation layer).
    \item A safety-bounded SON translation layer that clips per-update
    CIO changes to $\pm 3$~dB and the cumulative CIO to $\pm 9$~dB.
    \item A unit-test suite, smoke training run and diagnostic training
    run; the long multi-scenario run is currently executing on Google
    Colab Pro.
    \item A unified CSV ingestion schema that can accept both
    simulator output and Network Cell Info Lite drive-test logs.
\end{itemize}

\section{Why we need NTC's help}

Our framework is at the point where simulation-only results are no
longer enough to claim practical value. To demonstrate that the model
predicts and optimises handovers better than conventional A3+TTT
policies in the real Pokhara environment, we need access to a small
amount of carefully scoped operator-side data. We have come to NTC
because we believe NTC's continued support for academic research, and
its strong technical presence in Pokhara, make it the natural partner
for this validation.

\section{What we are asking for (specific items)}

\begin{enumerate}[leftmargin=*]
    \item Anonymised handover-event logs (no subscriber identifiers)
    for a representative time window over a small number of sites
    chosen jointly by NTC and our team, covering:
    \begin{itemize}[leftmargin=*]
        \item the urban core (Lakeside / Mahendrapul / Prithvi Chowk),
        \item the Prithvi Highway corridor outside Pokhara,
        \item the hilly cluster (Bindabasini / Sarangkot / Begnas).
    \end{itemize}
    \item Per-cell counters at the granularity NTC considers
    appropriate (for example, hourly summaries): handover success
    rate, ping-pong rate, RLF rate, PRB utilization and connected-UE
    count.
    \item The current A3 offset, TTT and CIO configuration table for
    the selected sites.
    \item Permission and any safety guidance for us to perform passive
    drive-test logging with a consumer Android handset on the same
    routes (no active calls, no SIM-card cloning, only reading what the
    phone already exposes).
\end{enumerate}

\section{How we will use and protect the data}

\begin{itemize}[leftmargin=*]
    \item Used \textbf{only} for the academic major project; never sold,
    shared with third parties, or used commercially.
    \item Stored encrypted on a private repository accessible only to
    the project team and the supervisor.
    \item Reported in the final document and in the defence presentation
    \textbf{only as aggregated, anonymised statistics}, never as raw
    logs, raw PCI lists or raw configuration tables.
    \item Deleted within $30$ days of project defence, on NTC's written
    confirmation.
    \item Subject to a Non-Disclosure Agreement on terms acceptable to
    NTC; the supervisor will co-sign.
\end{itemize}

\section{What NTC gets in return}

\begin{itemize}[leftmargin=*]
    \item A free, written technical report tailored to the Pokhara
    cluster, including before/after analysis of handover anomalies and
    suggested CIO/TTT adjustment ranges, with all of the safety bounds
    we have built into the SON layer respected by design.
    \item A presentation to the NTC Pokhara optimisation team at the
    end of the project.
    \item A clean Python reference implementation of SON-GNN-DQN that
    NTC engineers can inspect or run on their own data, free of cost.
    \item First right of access to any follow-up O-RAN xApp work.
\end{itemize}

\section{Questions we would like to discuss in the meeting}

Below is the structured list of questions we would like to discuss with
the NTC LTE optimisation team. We have grouped them for efficient
discussion. Items marked with a star ($\star$) are the most important
for validating our simulator assumptions; items marked with a dagger
($\dagger$) are the most important for the data request.

\subsection{RAN and Core}

\begin{enumerate}[leftmargin=*]
    \item ($\star$) Which RAN vendors are currently deployed in the
    Pokhara cluster (Huawei, ZTE, Nokia, Ericsson)? Mixed or single
    vendor?
    \item ($\star$) What is the architecture of the core (flat EPC,
    regional MMEs / SGW-LBOs)?
    \item Are X2-based handovers used cluster-internally? Where do
    S1-based handovers come into play?
    \item ($\star$) What is the typical inter-site distance in (a) the
    urban core, (b) the Prithvi Highway corridor and (c) the hilly
    cluster?
\end{enumerate}

\subsection{Handover Configuration}

\begin{enumerate}[leftmargin=*,resume]
    \item ($\star$, $\dagger$) What are the currently configured
    Event~A3 offset and Time-to-Trigger (TTT) values in the Pokhara
    cluster? Cluster-wide or per-cell?
    \item ($\star$, $\dagger$) Is the Cell Individual Offset (CIO) used
    per neighbour? In what range?
    \item Which thresholds are used for Event~A2?
    \item What hysteresis values are configured?
    \item Are different settings used for highway versus urban versus
    hilly clusters? How often are they re-tuned?
\end{enumerate}

\subsection{Mobility Robustness and Load Balancing}

\begin{enumerate}[leftmargin=*,resume]
    \item Is Mobility Robustness Optimization (MRO) enabled? Manual or
    automatic?
    \item Is Mobility Load Balancing (MLB) enabled? Manual or
    automatic? Which counters drive the decisions?
    \item ($\star$) What ping-pong and RLF rates are typically seen in
    the urban core during peak hours?
    \item Is sticky-cell or PCI confusion a known issue in any specific
    Pokhara sector?
    \item How often does the optimisation team re-tune the cluster?
\end{enumerate}

\subsection{KPIs and OSS Counters}

\begin{enumerate}[leftmargin=*,resume]
    \item ($\dagger$) Which KPIs and counters are tracked at the OSS
    level (handover success rate, ping-pong rate, RLF, PRB utilization,
    connected-UE count, throughput per cell)?
    \item What is the typical granularity of these counters
    (per cell / per minute / per hour)?
    \item ($\dagger$) Which of these counters could be shared, in
    anonymised form, for academic use under a clear non-disclosure
    scope?
\end{enumerate}

\subsection{Voice and Coverage}

\begin{enumerate}[leftmargin=*,resume]
    \item What is the approximate share of VoLTE versus CSFB in the
    Pokhara cluster?
    \item Where are the known coverage holes or weak-signal pockets
    (urban / highway / hilly)?
\end{enumerate}

\subsection{AI/ML and O-RAN Roadmap}

\begin{enumerate}[leftmargin=*,resume]
    \item Is any AI- or ML-based optimization currently in production
    in NTC's LTE network?
    \item Is NTC evaluating O-RAN-style disaggregated RAN or Near-RT
    RIC technology? What is the indicative timeline?
    \item Would NTC be interested in piloting an academic xApp on a
    test cell?
\end{enumerate}

\subsection{Data Sharing for the Project}

\begin{enumerate}[leftmargin=*,resume]
    \item ($\dagger$) Would NTC be willing to share, for a clearly
    defined academic window, (a) anonymised handover-event logs,
    (b) per-cell PRB-utilization counters and (c) the A3 / TTT / CIO
    parameter table for a small set of sites in Pokhara?
    \item ($\dagger$) Under what NDA and data-handling terms would such
    sharing be possible?
    \item ($\dagger$) Who in the Pokhara cluster would be the point of
    contact for the data exchange and the signoff?
\end{enumerate}

\section{Suggested next steps after the meeting}

\begin{enumerate}[leftmargin=*]
    \item NTC indicates which data items in Section~5 can be shared and
    under what terms.
    \item Both sides agree on a Non-Disclosure Agreement template.
    \item The NTC point of contact and our team agree on a delivery
    format and channel (encrypted ZIP, secure share link, on-site
    visit).
    \item Our team performs the analysis and shares the draft technical
    report with NTC for review before any public defence.
\end{enumerate}

\vspace{0.8cm}
\noindent Thank you for your time and continued support for academic
research in Nepal.

\vspace{0.8cm}
\noindent\begin{tabular}{ll}
Project Team: & Rojin Puri (PAS078BEI028)\\
              & Rujan Subedi (PAS078BEI030)\\
              & Saange Tamang (PAS078BEI031)\\
              & Sulav Kandel (PAS078BEI041)\\[0.4cm]
Supervisor:   & Asst.\ Prof.\ Er.\ Hom Nath Tiwari\\
              & Head of Department, Electronics and Computer Engineering\\
              & Paschimanchal Campus, IOE, Tribhuvan University\\
\end{tabular}

\end{document}
